\chapter{Yazılım Haritası: Python, SPSS ve R}
\label{app:software-map}

Bu ek, bölüm sonlarına eklenen R ile çözümlü uygulamaları tamamlayan kısa bir eşleştirme haritasıdır. Her bölüm kendi R kodunu, kontrol değerlerini ve yorumunu içerir; buradaki tablo temel analizlere karşılık gelen fonksiyon ailelerini topluca gösterir.
Kitapta kullanılan başlıca açık kaynak altyapılar NumPy, SciPy, pandas,
Matplotlib ve statsmodels'dır \citep{harris2020,virtanen2020,mckinney2022,hunter2007,seabold2010};
R tarafındaki veri düzenleme yaklaşımı için ayrıca bkz. \citet{wickham2017}.

\begin{table}[H]
\centering
\scriptsize
\renewcommand{\arraystretch}{1.08}
\begin{tabularx}{\textwidth}{@{}>{\raggedright\arraybackslash}p{0.16\textwidth}>{\raggedright\arraybackslash}X>{\raggedright\arraybackslash}X>{\raggedright\arraybackslash}X@{}}
\toprule
\textbf{Konu} & \textbf{Python} & \textbf{SPSS yolu} & \textbf{R eşdeğeri} \\
\midrule
Tanımlayıcı istatistikler & \texttt{pandas}, \texttt{numpy} & Analyze $>$ Descriptive Statistics & \texttt{summary}, \texttt{mean}, \texttt{sd} \\
Eksik veri denetimi & \texttt{isna}, \texttt{MICEData} & Missing Value Analysis / Multiple Imputation & \texttt{is.na}, \texttt{mice} \\
Örnekleme benzetimi & \texttt{numpy.random} & Transform / benzetim sözdizimi & \texttt{sample}, \texttt{rnorm}, \texttt{replicate} \\
\ifimoders\else
Bootstrap & \texttt{numpy.random.choice} & Bootstrap veya yeniden örnekleme eklentileri & \texttt{boot}, \texttt{replicate} \\
Rastgeleleştirme testi & \texttt{numpy.random} & Randomization/Exact Tests & \texttt{sample}, \texttt{coin} \\
\fi

Ortalama için güven aralığı & \texttt{scipy.stats.t} & Explore / One-Sample T Test & \texttt{t.test} \\
Tek örneklem ve iki örneklem testleri & \texttt{ttest\_1samp}, \texttt{ttest\_ind} & Compare Means & \texttt{t.test} \\
Doğrusal regresyon & \texttt{linregress} & Analyze $>$ Regression $>$ Linear & \texttt{lm}, \texttt{summary.lm} \\
\ifimoders\else
Çoklu regresyon & \texttt{statsmodels.OLS} & Analyze $>$ Regression $>$ Linear & \texttt{lm}, \texttt{car::vif} \\
Tek yönlü ANOVA & \texttt{f\_oneway} & Analyze $>$ Compare Means $>$ One-Way ANOVA & \texttt{aov} \\
Welch ANOVA & \texttt{anova\_generic} & One-Way ANOVA $>$ Welch & \texttt{oneway.test} \\
Çoklu karşılaştırma & \texttt{tukey\_hsd} & One-Way ANOVA $>$ Post Hoc & \texttt{TukeyHSD}, \texttt{pairwise.t.test} \\
\fi

Ki-kare testi & \texttt{scipy.stats} & Crosstabs & \texttt{chisq.test} \\
Fisher kesin testi & \texttt{scipy.stats} & Exact Tests & \texttt{fisher.test} \\
Korelasyon & \texttt{pearsonr} & Analyze $>$ Correlate $>$ Bivariate & \texttt{cor.test} \\
Spearman korelasyonu & \texttt{spearmanr} & Bivariate Correlations $>$ Spearman & \texttt{cor.test} (Spearman yöntemi) \\
Mann--Whitney / Wilcoxon & \texttt{mannwhitneyu}, \texttt{wilcoxon} & Nonparametric Tests & \texttt{wilcox.test} \\
Kruskal--Wallis & \texttt{kruskal} & Independent Samples $>$ Kruskal--Wallis & \texttt{kruskal.test} \\
Cronbach alfa & \texttt{numpy}/\texttt{pandas} ile formül & Analyze $>$ Scale $>$ Reliability Analysis & temel R formülü, \texttt{psych::alpha} \\
Omega & faktör modeli & Ek yazılım gerekir & \texttt{psych::omega} \\
\bottomrule
\end{tabularx}
\end{table}

\section*{Asgari R Raporlama Şablonu}

Tekrarlanabilir bir R analizi; veri kaynağını, temizleme kararlarını, paket sürümlerini, rassal tohum değerini, model formülünü, tanı kontrollerini ve tablo ile şekilleri üreten tam kodu kaydetmelidir. Asgari bir proje betiği çoğu zaman şöyle başlar:

\begin{lstlisting}[style=prismR]
set.seed(2026)
sessionInfo()
veri <- data.frame(saat = c(2, 3, 4, 5, 6),
                   puan = c(58, 62, 65, 69, 70))
model <- lm(puan ~ saat, data = veri)
summary(model)
confint(model)
plot(model)
\end{lstlisting}

Bu şablondaki beş satır öğretim amaçlıdır; dış dosya gerekmez. Bölümlerdeki
R ile çözümlü öğretim uygulamaları kod içindeki girdilerle bağımsız çalışır; ortak gerçek veri uygulamaları ise eşlikçi CSV dosyasını okur. Haritadaki ek paket seçenekleri zorunlu değildir; bölüm çözümlerinde standart kurulumdaki işlevler kullanılır.

\section*{Çeviri İlkesi}

Python, SPSS ve R arasında geçiş yaparken önce araştırma sorusunu ve hedef
büyüklüğü, sonra veri yapısını, en son yazılım fonksiyonunu eşleştirin. Benzer
menü adları aynı varsayılanların kullanıldığı anlamına gelmez: eksik veri
işlemleri, eşit varyans seçimi, süreklilik düzeltmesi ve kesin test seçenekleri
çıktılar karşılaştırılmadan önce denetlenmelidir.

\section*{Eşlikçi paketi yeniden üretme}

Kitabın gerçek veri örnekleri \texttt{companion/} klasöründedir. Temiz veri,
sayısal tablolar, örnekleme benzetimleri ve grafikler şu komutla yeniden
oluşturulur:
\begin{center}
\ifimoders
\texttt{python companion/code/run\_imo301.py}
\else
\texttt{python companion/code/run\_all.py}
\fi
\end{center}
Çalışma sonunda yazılım sürümleri \texttt{environment.txt} dosyasına
kaydedilir. Üretilen dosyaların bütünlük özetleri ise
\texttt{manifest\_sha256.csv} içinde tutulur. Bir analiz sonucu paylaşılırken
kullanılan kitap baskısı, veri paketi sürümü ve kaynak veri birlikte
belirtilmelidir.